\chapter{Kalkulus Diferensial}\label{ch:b2:diffcalc}

Kalkulus dua peubah pada jilid Tahun ke-1 mematang menjadi
kalkulus \omterm{def:b2:diffcalc:differential}{diferensial} bagi pemetaan antarruang bernorma: dengan
\emph{\omterm{def:b2:diffcalc:differential}{diferensial}} sebagai hampiran linear terbaiknya, aturan rantai
dalam keumuman penuh, teorema kesimetrikan Schwarz yang \emph{dibuktikan},
rumus Taylor, dan analisis orde dua yang \omterm{def:b2:metric:complete}{lengkap} bagi ekstremumnya. Adapun
teorema fungsi balikan, yakni mahkota teorinya, dinyatakan dengan
strategi buktinya --- sebuah titik tetap Banach --- yang dibuat gamblang.

Di sepanjang bab ini, $U$ adalah himpunan bagian \omterm{def:b2:metric:topology}{terbuka} $\R^n$ (atau sebuah ruang
bernorma; sebab kasus berdimensi hingganya mengusung semua gagasannya), dan $f \colon U \to
\R^m$.

\section{Diferensialnya}

\begin{definition}\label{def:b2:diffcalc:differential}
Fungsi $f$ disebut \emph{terdiferensialkan}\index{terdiferensialkan} di $a$ apabila ada
pemetaan linear (yang \omterm{def:b2:metric:continuity}{kontinu}) $\dd f_a \colon \R^n \to \R^m$ dengan
\[
f(a + h) = f(a) + \dd f_a(h) + o(\norm h)
\qquad (h \to 0).
\]
Pemetaan $\dd f_a$, yakni \emph{diferensial}\index{diferensial}
$f$ di $a$, bersifat tunggal; dan matriksnya dalam basis kanoniknya adalah
\emph{matriks Jacobi}\index{matriks Jacobi} $J_f(a) =
\bigl(\frac{\partial f_i}{\partial x_j}(a)\bigr)$.
Keterdiferensialan mengakibatkan kekontinuan beserta keberadaan semua
turunan berarah $\dd f_a(v) = \lim_{t\to0}\frac{f(a + tv)
- f(a)}{t}$; sedangkan konversnya gagal
(\cref{exo:b2:diffcalc:2}). Untuk $m = 1$, berlaku $\dd f_a(h) = \langle
\nabla f(a), h\rangle$: yakni gradien Tahun ke-1, yang kini dipahami sebagai
vektor yang mewakili diferensialnya.
\end{definition}

\begin{theorem}[Kriteria $C^1$]\label{thm:b2:diffcalc:c1}
Bila semua turunan parsial $f$ ada pada $U$ dan \omterm{def:b2:metric:continuity}{kontinu}
di $a$, maka $f$ \omterm{def:b2:diffcalc:differential}{terdiferensialkan} di $a$. Jadi ``$C^1$ pada $U$'' ---
yakni turunan parsial yang \omterm{def:b2:metric:continuity}{kontinu} --- mengakibatkan keterdiferensialan
di mana-mana, dengan \omterm{def:b2:diffcalc:differential}{diferensial} yang \omterm{def:b2:metric:continuity}{kontinu}.
\end{theorem}

\begin{proof}
Komponen demi komponen (jadi $m = 1$ sudah cukup). Bukti Tahun ke-1 untuk dua
peubah --- yakni gerakkan satu koordinat sekali waktu, terapkan teorema nilai
rerata satu peubah pada tiap kakinya, lalu pakai kekontinuan turunan parsialnya
di $a$ --- menyamarata kata demi kata ke $n$ kaki:
\[
f(a + h) - f(a) = \sum_{j=1}^{n} \bigl(f(a + h^{(j)}) - f(a +
h^{(j-1)})\bigr)
= \sum_j h_j\,\frac{\partial f}{\partial x_j}(\xi_j),
\]
dengan $h^{(j)}$ membekukan $j$ koordinat pertama $h$ dan
$\xi_j$ terletak pada kaki ke-$j$-nya; lalu kekontinuannya mengubah tiap
$\frac{\partial f}{\partial x_j}(\xi_j)$ menjadi $\frac{\partial
f}{\partial x_j}(a) + o(1)$, dan galatnya bernilai $o(\norm h)$.
\end{proof}

\begin{theorem}[Aturan rantai]\label{thm:b2:diffcalc:chainrule}
Bila $f$ \omterm{def:b2:diffcalc:differential}{terdiferensialkan} di $a$ dan $g$ di $f(a)$, maka $g \circ f$
\omterm{def:b2:diffcalc:differential}{terdiferensialkan} di $a$ dengan
\[
\dd(g \circ f)_a = \dd g_{f(a)} \circ \dd f_a ,
\qquad
J_{g\circ f}(a) = J_g\bigl(f(a)\bigr)\,J_f(a) :
\]
jadi \omterm{def:b2:diffcalc:differential}{matriks Jacobinya} mengalikan.\index{aturan rantai}
\end{theorem}

\begin{proof}
Tulislah $f(a + h) = f(a) + \dd f_a(h) + \norm h\,\varepsilon_1(h)$
dan $g(b + k) = g(b) + \dd g_b(k) + \norm k\,\varepsilon_2(k)$ dengan
$b = f(a)$ dan $k = k(h) = \dd f_a(h) + \norm h \varepsilon_1(h)$.
Setelah disulihkan,
\[
g(f(a+h)) = g(b) + \dd g_b\bigl(\dd f_a(h)\bigr)
+ \norm h\,\dd g_b(\varepsilon_1(h)) + \norm{k}\,\varepsilon_2(k),
\]
dan kedua suku galatnya bernilai $o(\norm h)$: yang pertama karena $\dd g_b$
\omterm{def:b2:metric:continuity}{kontinu} dan $\varepsilon_1 \to 0$; sedangkan yang kedua karena $\norm k
\leq C\norm h$ (yakni pemetaan linear terbatas ditambah suku kecil) dan
$\varepsilon_2(k) \to 0$ ketika $h \to 0$.
\end{proof}

\begin{example}[Fungsi radial, sekali untuk selamanya]\label{ex:b2:diffcalc:radial}
Misalkan $r(x) = \norm x_2$ pada $\R^n\setminus\{0\}$ dan $f = g \circ
r$ dengan $g$ fungsi $C^1$ satu peubah. Pertama, $r$ bersifat
\omterm{def:b2:diffcalc:differential}{terdiferensialkan} jauh dari $0$: sebab dari $r^2 = \sum x_i^2$,
\[
\frac{\partial r}{\partial x_i} = \frac{x_i}{r},
\qquad\text{yakni}\qquad
\nabla r(x) = \frac{x}{\norm x} ,
\]
yakni vektor radial satuannya (turunkan $r^2$ lalu bagilah --- atau
terapkan aturan rantai pada $\sqrt{\cdot}$). Lalu aturan rantainya
memberi, untuk setiap fungsi radial,
\[
\nabla f(x) = g'\bigl(\norm x\bigr)\,\frac{x}{\norm x} .
\]
Contoh yang dikerjakan: $g(r) = \frac1r$ menghasilkan $\nabla\frac{1}{\norm
x} = -\frac{x}{\norm x^3}$, yakni medan kuadrat-terbalik milik
gravitasi dan elektrostatika --- dengan arah radial dan besar
$\frac{1}{\norm x^2}$. Pelajaran penutupnya: gradien fungsi
radial bersifat radial karena himpunan arasnya berupa bola dan
gradiennya ortogonal terhadap himpunan aras; sedangkan di $x = 0$,
sebaliknya, $r$ \emph{tidak} \omterm{def:b2:diffcalc:differential}{terdiferensialkan} (sebab tak ada calon
pemetaan linear yang cocok dengan $\norm h$ dari segala arah) --- jadi
profil radial yang mulus menuntut $g'(0) = 0$ agar dapat menyeberangi titik asalnya
dengan anggun.
\end{example}

\begin{theorem}[Ketaksamaan nilai rerata]\label{thm:b2:diffcalc:mvi}
Misalkan $f$ \omterm{def:b2:diffcalc:differential}{terdiferensialkan} pada $U$ dan misalkan ruas $\intcc{a}{b}
= \{a + t(b-a)\}$ terletak di $U$. Maka
\[
\norm{f(b) - f(a)} \leq \norm{b - a}\,
\sup_{x \in \intcc{a}{b}} \vertiii{\dd f_x} .
\]
Khususnya, pemetaan \omterm{def:b2:diffcalc:differential}{terdiferensialkan} yang \omterm{def:b2:diffcalc:differential}{diferensialnya} nol pada himpunan
\omterm{def:b2:metric:topology}{terbuka} yang \emph{\omterm{def:b2:metric:connected}{terhubung}} bersifat tetap.
\end{theorem}

\begin{proof}
Fungsi $\varphi(t) = f(a + t(b-a))$ \omterm{def:b2:diffcalc:differential}{terdiferensialkan} pada
$\intcc{0}{1}$ dengan $\varphi'(t) = \dd f_{a + t(b-a)}(b - a)$
(lewat aturan rantai), yang bernorma $\leq M\norm{b-a}$ dengan $M$ menyatakan
sup pada displainya. Untuk nilai di $\R$, ketaksamaan nilai rerata satu peubah bagi $f$
merampungkannya; sedangkan untuk nilai vektornya terapkanlah ia pada $t \mapsto \langle u,
\varphi(t)\rangle$ dengan $u$ vektor satuan sepanjang $f(b) - f(a)$.
Ketetapannya: sebab ia tetap secara lokal (lewat ruas di dalam bola) ditambah keterhubungannya
(sebab himpunan tempat $f$ bernilai tertentu bersifat \omterm{def:b2:metric:topology}{terbuka} sekaligus tertutup:
\cref{ch:b2:metric}).
\end{proof}

\begin{example}[Sebuah konstanta Lipschitz dari ketaksamaan nilai rerata]\label{ex:b2:diffcalc:lipschitz}
Apakah $f(x, y) = \sin x\,\sin y$ bersifat \omterm{def:b2:metric:continuity}{Lipschitz} pada $\R^2$, dan dengan
konstanta yang mana? Gradiennya $\nabla f = (\cos x\sin y,\
\sin x\cos y)$, yang kuadrat \omterm{def:b2:nvs:norm}{normanya}
\[
\cos^2x\sin^2y + \sin^2x\cos^2y
\leq \sin^2 y + \cos^2y\cdot 1 = 1
\]
(sebab $\cos^2x$ dan $\sin^2x$ dibatasi $1$ secara terpisah), jadi
$\vertiii{\dd f_{(x,y)}} = \norm{\nabla f} \leq 1$ di mana-mana,
lalu \cref{thm:b2:diffcalc:mvi} pada ruas antara sembarang dua
titik memberi
\[
\abs{f(b) - f(a)} \leq \norm{b - a}_2 :
\]
jadi $f$ bersifat \omterm{def:b2:metric:continuity}{Lipschitz} berkonstanta $1$, dan konstantanya tajam (sebab di dekat
titik asalnya, $f(x, \tfrac\pi2) = \sin x$ berlereng $1$). Pelajaran
penutupnya: ketaksamaan nilai rerata mengubah batas \omterm{def:b2:funcseq:def}{titik demi titik}
pada \omterm{def:b2:diffcalc:differential}{diferensialnya} menjadi modulus kekontinuan yang global ---
yakni rute baku menuju taksiran \omterm{def:b2:metric:continuity}{Lipschitz} pada setiap dimensi,
sekaligus mesin di dalam \cref{exo:b2:diffcalc:12}.
\end{example}

\section{Turunan kedua}

\begin{theorem}[Schwarz]\label{thm:b2:diffcalc:schwarz}
Bila $f$ bersifat $C^2$ pada $U$ (yakni semua turunan parsial keduanya ada dan
\omterm{def:b2:metric:continuity}{kontinu}), maka untuk setiap $i, j$:\index{teorema Schwarz}
\[
\frac{\partial^2 f}{\partial x_i\,\partial x_j}
= \frac{\partial^2 f}{\partial x_j\,\partial x_i} .
\]
\end{theorem}

\begin{proof}
Dua peubah sudah cukup (dengan $x = x_i$, $y = x_j$, dan sisanya dibekukan).
Tinjaulah selisih keduanya
\[
\Delta(h) = f(a + h, b + h) - f(a + h, b) - f(a, b + h) + f(a,b) .
\]
Tetapkan $h$ lalu tetapkan $\varphi(x) = f(x, b+h) - f(x, b)$: maka $\Delta(h)
= \varphi(a + h) - \varphi(a)$, dan dua kali penerapan teorema nilai
rerata memberi
\[
\Delta(h) = h\,\varphi'(\xi)
= h\Bigl(\frac{\partial f}{\partial x}(\xi, b+h) -
\frac{\partial f}{\partial x}(\xi, b)\Bigr)
= h^2\,\frac{\partial^2 f}{\partial y\,\partial x}(\xi, \eta),
\]
dengan $\xi \in \intoo{a}{a+h}$ dan $\eta \in \intoo{b}{b+h}$. Berkat
kekontinuannya, $\frac{\Delta(h)}{h^2} \to \frac{\partial^2 f}{\partial
y\partial x}(a, b)$ ketika $h \to 0$. Perhitungan yang sama dengan
peran peubahnya dipertukarkan (bekukan peubah keduanya lebih dulu)
memberi $\frac{\Delta(h)}{h^2} \to \frac{\partial^2 f}{\partial x
\partial y}(a,b)$: jadi kedua limit besaran yang sama itu berimpit.
\end{proof}

\begin{example}[Mengapa $C^2$ diperlukan: contoh tandingan Peano]\label{ex:b2:diffcalc:peano}
Misalkan $f(x, y) = \dfrac{xy(x^2 - y^2)}{x^2 + y^2}$ dengan $f(0,0) = 0$.
Jauh dari titik asalnya $f$ bersifat $C^\infty$; sedangkan di titik asalnya semua turunan
parsial pertama dan keduanya ada, tetapi yang campurannya tak sepakat. Hitunglah
sepanjang sumbunya: $f(x, 0) = f(0, y) = 0$, dan untuk $y \neq 0$,
\[
\frac{\partial f}{\partial x}(0, y)
= \lim_{x\to0}\frac{f(x,y)}{x}
= \frac{y(0 - y^2)}{y^2} = -y ,
\qquad\text{secara simetris}\qquad
\frac{\partial f}{\partial y}(x, 0) = x .
\]
Karena itu
\[
\frac{\partial^2 f}{\partial y\,\partial x}(0,0)
= \frac{\dd}{\dd y}\Bigl[\frac{\partial f}{\partial
x}(0,y)\Bigr]_{y=0} = -1,
\qquad
\frac{\partial^2 f}{\partial x\,\partial y}(0,0) = +1 :
\]
jadi kedua turunan parsial campurannya ada dan berbeda. Dan tak ada pertentangan dengan
\cref{thm:b2:diffcalc:schwarz}: sebab turunan parsial kedua $f$
tidak \omterm{def:b2:metric:continuity}{kontinu} di $0$ (ujilah sepanjang $y = tx$). Pelajaran penutupnya:
teorema Schwarz adalah teorema sejati tentang \emph{kekontinuan},
bukan kesamaan formal --- dan hipotesis ``$C^2$'' pada
Taylor--Young di bawah sungguh bekerja.
\end{example}

\begin{theorem}[Taylor--Young pada orde 2]\label{thm:b2:diffcalc:taylor}
Misalkan $f \colon U \to \R$ bersifat $C^2$ dan $a \in U$. Maka, ketika $h \to 0$,
\[
f(a + h) = f(a) + \langle\nabla f(a), h\rangle
+ \frac12\, \langle H_a h,\, h\rangle + o\bigl(\norm h^2\bigr),
\]
dengan $H_a = \bigl(\frac{\partial^2 f}{\partial x_i\partial
x_j}(a)\bigr)$ menyatakan \emph{matriks Hessian}\index{matriks Hessian}
(yang \omterm{def:b2:quadratic:adjoint}{simetrik}, berkat Schwarz).
\end{theorem}

\begin{proof}
Terapkan teorema Taylor--Young satu peubah (jilid Tahun ke-1) pada
$\varphi(t) = f(a + th)$ di $\intcc{0}{1}$: sebab menurut aturan rantai,
$\varphi'(t) = \langle \nabla f(a + th), h\rangle$ dan $\varphi''(t)
= \langle H_{a+th}h, h\rangle$, yang keduanya \omterm{def:b2:metric:continuity}{kontinu} dalam $t$. Lalu
$\varphi(1) = \varphi(0) + \varphi'(0) + \frac12\varphi''(\theta)$
(lewat Taylor--Lagrange) dengan $\theta \in \intoo{0}{1}$, dan kekontinuan
turunan parsial keduanya mengubah $\varphi''(\theta) =
\varphi''(0) + o(1)\cdot\norm h^2$ secara \omterm{def:b2:funcseq:def}{seragam}: jadi diperoleh
uraian pada displainya.
\end{proof}

\begin{example}[Sebuah uraian lewat dua cara]\label{ex:b2:diffcalc:taylortwoways}
Uraikan $f(x, y) = \eu^x\cos y$ di titik asalnya sampai orde $2$.
\emph{Lewat komposisi uraian satu peubah:}
\[
\eu^x\cos y = \Bigl(1 + x + \frac{x^2}{2} +
o(x^2)\Bigr)\Bigl(1 - \frac{y^2}{2} + o(y^2)\Bigr)
= 1 + x + \frac{x^2 - y^2}{2} + o\bigl(\norm{(x,y)}^2\bigr) .
\]
\emph{Lewat turunan parsial:} $f_x = \eu^x\cos y$ dan $f_y =
-\eu^x\sin y$, jadi $\nabla f(0) = (1, 0)$; sedangkan $f_{xx} = f$,
$f_{yy} = -f$, $f_{xy} = -\eu^x\sin y$ memberi $H_0 =
\operatorname{diag}(1, -1)$: jadi \cref{thm:b2:diffcalc:taylor}
menghasilkan kembali $1 + x + \frac12(x^2 - y^2)$. Kedua perhitungannya
sepakat, dan rute komposisinya lebih cepat --- tanpa satu turunan
parsial kedua pun. Pelajaran penutupnya: titik asalnya \emph{bukan}
titik kritis (sebab $\nabla f \neq 0$), jadi kendati Hessiannya tak tentu
tak ada pelana yang dapat diumumkan: sebab suku linearnya yang berkuasa,
dan uji orde duanya hanya pernah berbicara di titik
kritis.
\end{example}

\begin{theorem}[Uji ekstremum orde dua, dibuktikan]\label{thm:b2:diffcalc:extrema}
Misalkan $f$ bersifat $C^2$ di dekat titik kritis $a$ (dengan $\nabla f(a) = 0$),
dan berhessian $H = H_a$.
\begin{enumerate}
  \item Bila $H$ definit positif, maka $a$ merupakan minimum lokal yang tegas
        (dan bila definit negatif: maksimum).
  \item Bila $H$ punya \omterm{def:b2:reduction:eigen}{nilai eigen} bertanda keduanya, maka $a$ berupa pelana: jadi tanpa
        ekstremum.
  \item Bila $H$ singular (dan semidefinit), maka tak ada kesimpulan.
\end{enumerate}
Adapun uji ``$rt - s^2$'' Tahun ke-1 adalah kasus $n = 2$: sebab $\det H = rt -
s^2$, dan tanda $\operatorname{tr}$-nya dibaca dari $r$.
\end{theorem}

\begin{proof}
Menurut teorema spektral (\cref{thm:b2:quadratic:spectral}), \omterm{def:b2:quadratic:def}{bentuk
kuadratik} $H$ terjepit di antara \omterm{def:b2:reduction:eigen}{nilai eigen} ekstremnya:
$\lambda_{\min}\norm h^2 \leq \langle Hh, h\rangle \leq
\lambda_{\max}\norm h^2$.

(1) Bila $\lambda_{\min} > 0$: maka Taylor--Young memberi
\[
f(a + h) - f(a) \geq \frac{\lambda_{\min}}{2}\norm h^2 -
o(\norm h^2) > 0
\]
untuk $h \neq 0$ yang kecil: jadi minimum lokal yang tegas.

(2) Sepanjang \omterm{def:b2:reduction:eigen}{vektor eigen} $v_+$ dengan $\lambda_+ > 0$: berlaku $f(a + tv_+) -
f(a) = \frac{\lambda_+}{2}t^2 + o(t^2) > 0$ untuk $t$ yang kecil; sedangkan sepanjang
$v_-$ dengan $\lambda_- < 0$ selisihnya negatif: jadi kedua tandanya
muncul pada setiap persekitarannya.

(3) Fungsi $f(x,y) = x^2 + y^4$, $x^2 - y^4$, $x^2 + y^3$ berbagi Hessian
singular semidefinit yang sama di $0$ dengan tiga kelakuan yang
berbeda.
\end{proof}

\begin{example}[Sebuah penggolongan lengkap, termasuk yang global]\label{ex:b2:diffcalc:quarticrun}
Golongkanlah semua ekstremum $f(x, y) = x^4 + y^4 - 4xy$ pada $\R^2$.
\emph{Titik kritisnya:} $\nabla f = (4x^3 - 4y,\ 4y^3 - 4x) =
0$ memberi $y = x^3$ dan $x = y^3 = x^9$, jadi $x(x^8 - 1) = 0$:
sehingga penyelesaian realnya adalah $(0,0)$, $(1,1)$, $(-1,-1)$.
\emph{Hessiannya:} $H = \begin{pmatrix} 12x^2 & -4\\ -4 &
12y^2\end{pmatrix}$. Di $(\pm1, \pm1)$: $\begin{pmatrix} 12 &
-4\\ -4 & 12\end{pmatrix}$, bernilai eigen $8$ dan $16$: jadi
definit positif, sehingga minimum lokal yang tegas dengan $f = -2$. Di $(0,0)$:
$\begin{pmatrix} 0 & -4\\ -4 & 0\end{pmatrix}$, bernilai eigen
$\pm4$: jadi sebuah pelana. \emph{Kegloballannya:} dari $2\abs{xy} \leq x^2 +
y^2$,
\[
f(x, y) \geq x^4 + y^4 - 2(x^2 + y^2)
= (x^2 - 1)^2 + (y^2 - 1)^2 + x^2 + y^2 - 2
\xrightarrow[\norm{(x,y)}\to\infty]{} +\infty :
\]
jadi $f$ bersifat koersif, sehingga ia mencapai minimum global
(berkat kekompakan himpunan subarasnya), yang niscaya di sebuah titik
kritis: jadi nilai $-2$, baik di $(1,1)$ maupun di $(-1,-1)$, adalah
minimum globalnya; dan tak ada maksimum (sebab $f$ tak terbatas ke atas).
Pelajaran penutupnya: uji lokalnya menggolongkan calonnya, tetapi
hanya hujah pertumbuhan yang mengubah ``lokal'' menjadi ``global'' --- yakni
pola dua langkah setiap bukti pengoptimalan pada buku ini.
\end{example}

\begin{method}[Menggolongkan ekstremum $f \colon \R^n \to \R$]\label{met:b2:diffcalc:extrema}
\begin{enumerate}
  \item Selesaikan $\nabla f = 0$ (yakni semua titik kritisnya; dan pada daerah
        yang berbatas, tanganilah perbatasannya secara terpisah seperti pada
        \cref{exo:b2:diffcalc:7}).
  \item Di tiap titik kritisnya, hitunglah Hessiannya beserta
        tanda \omterm{def:b2:reduction:eigen}{nilai eigennya} --- pada dimensi $2$, cukuplah $\det H$
        dan $\operatorname{tr} H$: bila $\det < 0$ maka pelana; bila $\det >
        0$ maka ekstremum, dengan jenis yang diberikan tanda
        tracenya; dan bila $\det = 0$ ujinya membisu, jadi kajilah $f$
        sepanjang kurva.
  \item Untuk pernyataan global, tambahkanlah hujah kekompakan atau kekoersifan
        (yakni $f \to +\infty$ di tak hingga, atau himpunan kendala yang
        \omterm{def:b2:metric:compact}{kompak}), lalu bandingkanlah nilai kritisnya.
\end{enumerate}
\end{method}

\section{Teorema fungsi balikan}

\begin{theorem}[Teorema fungsi balikan]\label{thm:b2:diffcalc:inverse}
Misalkan $f \colon U \to \R^n$ bersifat $C^1$ dan $a \in U$ dengan $\dd f_a$
yang \emph{terbalikkan}. Maka ada persekitaran \omterm{def:b2:metric:topology}{terbuka} $V \ni a$ dan $W
\ni f(a)$ sedemikian sehingga $f \colon V \to W$ merupakan bijeksi dengan
balikan $C^1$, dan\index{teorema fungsi balikan}
\[
\dd (f^{-1})_{f(x)} = (\dd f_x)^{-1} \qquad (x \in V).
\]
\end{theorem}
\admitted

\begin{remark}[Mengapa ia benar: strategi titik tetapnya]
Menyelesaikan $f(x) = y$ di dekat $a$ ditulis ulang sebagai persamaan titik tetap
$x = x + \dd f_a^{-1}\bigl(y - f(x)\bigr) =: \Phi_y(x)$; lalu pemetaan
$\Phi_y$ berdiferensial $\mathrm{id} - \dd f_a^{-1}\dd f_x$,
yang kecil di dekat $a$ berkat kekontinuan $\dd f$, jadi $\Phi_y$ merupakan
kontraksi pada bola tertutup yang kecil sehingga teorema titik tetap Banach
(\cref{thm:b2:metric:banach}) menyediakan penyelesaian lokal
tunggalnya $x = f^{-1}(y)$. Adapun kekontinuan dan keterdiferensialan
balikannya lalu menyusul dari taksiran di dalam kontraksinya. Tata buku
selengkapnya dikerjakan pada Tahun ke-3; sedangkan strateginya --- dan
pernyataannya --- dipakai secara bebas mulai kini. Adapun pendampingnya, yakni
\emph{teorema fungsi implisit} (yang menyelesaikan $F(x, y) = 0$ bagi $y(x)$
ketika $\frac{\partial F}{\partial y}$ terbalikkan), menyusul dengan
menerapkan teoremanya pada $(x, y) \mapsto (x, F(x,y))$.
\end{remark}

\begin{example}[Koordinat kutub]\label{ex:b2:diffcalc:polar}
Pemetaan $\Phi(r, \theta) = (r\cos\theta, r\sin\theta)$ berjacobi
\[
J_\Phi = \begin{pmatrix} \cos\theta & -r\sin\theta\\ \sin\theta &
r\cos\theta \end{pmatrix},
\qquad \det J_\Phi = r :
\]
yang terbalikkan untuk $r \neq 0$, jadi $\Phi$ merupakan difeomorfisma $C^1$
lokal jauh dari titik asalnya --- yakni izin untuk ``berpindah ke
koordinat kutub'', yang diperbarui bagi integral lipat pada
\cref{ch:b2:multint}.
\end{example}

\begin{example}[Lokal di mana-mana, global tak di mana pun]\label{ex:b2:diffcalc:localnotglobal}
Misalkan $f(x, y) = \bigl(\eu^x\cos y,\ \eu^x\sin y\bigr)$ pada
$\R^2$. \omterm{def:b2:diffcalc:differential}{Matriks Jacobinya},
\[
J_f = \begin{pmatrix}
\eu^x\cos y & -\eu^x\sin y\\
\eu^x\sin y & \eu^x\cos y
\end{pmatrix},
\qquad
\det J_f = \eu^{2x} > 0 ,
\]
tak pernah lenyap: jadi menurut \cref{thm:b2:diffcalc:inverse}, $f$ merupakan
difeomorfisma $C^1$ lokal di \emph{setiap} titik bidangnya.
Namun $f$ jauh dari injektif: sebab $f(x, y + 2\pi) = f(x, y)$, jadi
setiap nilainya diambil tak berhingga kali; dan ia juga tak surjektif,
sebab $\norm{f(x, y)} = \eu^x > 0$ meleset dari titik asalnya.
Pelajaran penutupnya: teorema fungsi balikan bersifat
\emph{lokal} tanpa dapat direduksi --- sebab keterbalikan setiap $\dd f_a$ menghasilkan
tambal sulam balikan lokal yang tak harus terangkai menjadi satu.
(Pembaca yang mengenal bilangan kompleks akan mengenali $z \mapsto
\eu^z$; dan tambal sulamnya adalah keluarga cabang logaritmanya.)
Bandingkan dengan \cref{exo:b2:diffcalc:12}, tempat sebuah hipotesis global
yang kuantitatif memang memaksa satu balikan global.
\end{example}

\begin{remark}[Jebakan yang sering muncul]
\emph{(i) Turunan berarah itu murah, sedangkan \omterm{def:b2:diffcalc:differential}{diferensial}
tidak:} semua turunan berarahnya boleh saja ada --- bahkan gagal
bergantung secara linear pada arahnya --- tanpa
keterdiferensialan (\cref{exo:b2:diffcalc:2}); jadi hanya kriteria $C^1$
(\cref{thm:b2:diffcalc:c1}) yang meningkatkan turunan parsial menjadi sebuah
\omterm{def:b2:diffcalc:differential}{diferensial}. \emph{(ii) Kritis tak berarti ekstremal:}
sebab pelana (\cref{ex:b2:diffcalc:quarticrun}) dan kasus singular
yang membisu (\cref{thm:b2:diffcalc:extrema}~(3)) sama-sama bersembunyi
di balik $\nabla f = 0$. \emph{(iii) Tak ada kesamaan nilai rerata
yang bernilai vektor:} sebab hanya \emph{ketaksamaan} pada
\cref{thm:b2:diffcalc:mvi} yang bertahan
(\cref{exo:b2:diffcalc:9}); jadi jangan pernah menulis $f(b) - f(a) = \dd
f_c(b-a)$ untuk $f$ yang bernilai di $\R^m$ dengan $m \geq 2$.
\emph{(iv) Keterbalikan lokal bukan keinjektifan:} lihat
\cref{ex:b2:diffcalc:localnotglobal}. \emph{(v) Gradiennya
milik hasil kali dalamnya:} sebab $\nabla f$ adalah vektor yang
mewakili $\dd f_a$ dalam hasil kali dalam yang dipilih; jadi ubahlah
hasil kalinya (seperti pada contoh berbobot bab \omterm{def:b2:quadratic:def}{bentuk kuadratik})
dan gradiennya berputar, sedangkan \omterm{def:b2:diffcalc:differential}{diferensialnya} ---
yakni objek yang hakiki --- tak bergerak.
\end{remark}

\begin{remark}[Di mana ini dipakai]
Segala yang mengalir dari bab ini adalah kalkulus \omterm{def:b2:diffcalc:differential}{diferensial}
yang diterapkan: bab persamaan \omterm{def:b2:diffcalc:differential}{diferensial} melinearkan aliran lalu
memakai rumus \omterm{def:b2:linalg:det}{determinan} Liouville (yang dibuktikan pada soal akhir pekan
bab ini); bab kurva dan permukaan mengkaji
himpunan aras dan parameterisasinya lewat teorema fungsi
implisit; dan integral lipat berganti peubah lewat \omterm{def:b2:diffcalc:differential}{matriks Jacobi}.
Adapun soal akhir pekan mengembangkan kalkulusnya \emph{pada ruang
matriks itu sendiri} --- yakni \omterm{def:b2:diffcalc:differential}{diferensial} \omterm{def:b2:linalg:det}{determinannya}, balikannya,
eksponensial matriksnya, dan grup ortogonal sebagai himpunan aras yang
mulus --- yakni bayangan Tahun ke-2 bagi apa yang diresmikan jilid
Tahun ke-3 sebagai manifold dan grup Lie.
\end{remark}

\section{Latihan}

\begin{exercise}[$\star$]\label{exo:b2:diffcalc:1}
Hitunglah \omterm{def:b2:diffcalc:differential}{matriks Jacobi} bagi $f(x,y) = (x^2 - y^2,\, 2xy)$ dan
bagi $\Phi(r,\theta,z) = (r\cos\theta, r\sin\theta, z)$; lalu di mana
\omterm{def:b2:diffcalc:differential}{diferensialnya} terbalikkan?
\end{exercise}

\begin{exercise}[$\star$]\label{exo:b2:diffcalc:2}
Misalkan $f(x,y) = \frac{x^3}{x^2 + y^2}$ dengan $f(0,0) = 0$. Buktikan bahwa semua
turunan berarah $f$ di $0$ ada, tetapi bahwa $f$ tidak
\omterm{def:b2:diffcalc:differential}{terdiferensialkan} di $0$ \emph{(sebab pemetaan $v \mapsto$ turunan
berarahnya tidak linear)}.
\end{exercise}

\begin{exercise}[$\star$]\label{exo:b2:diffcalc:3}
Carilah lalu golongkanlah titik kritis $f(x, y) = x^3 + y^3 -
3xy$ dengan memakai \cref{thm:b2:diffcalc:extrema}, dan juga bagi $g(x,y) = x^4 +
y^4 - 2(x - y)^2$.
\end{exercise}

\begin{exercise}[$\star\star$]\label{exo:b2:diffcalc:4}
Misalkan $f \colon \R^n \to \R$ bersifat $C^1$ dan \emph{homogen berderajat
$p$}: yakni $f(tx) = t^pf(x)$ untuk $t > 0$. Buktikan kesamaan Euler
\[
\langle \nabla f(x), x\rangle = p\,f(x) ,
\]
beserta konversnya bagi fungsi $C^1$ pada $\R^n\setminus\{0\}$.
\end{exercise}

\begin{exercise}[$\star\star$]\label{exo:b2:diffcalc:5}
Misalkan $A$ \omterm{def:b2:quadratic:adjoint}{simetrik} dan $f(x) = \frac12\langle Ax, x\rangle -
\langle b, x\rangle$. Hitunglah $\nabla f$ dan $H_f$; lalu kapan $f$
cembung? Dengan mengandaikan $A$ definit positif, tunjukkanlah bahwa $f$ punya satu
minimum global di penyelesaian $Ax = b$ --- yakni alasan keberadaan
penurunan gradien.
\end{exercise}

\begin{exercise}[$\star\star$]\label{exo:b2:diffcalc:6}
(Pengganda Lagrange, satu kendala, dibuktikan dengan tangan) Misalkan $f, g$
bersifat $C^1$ pada $\R^2$, dan andaikan $f$ mencapai, di $a$, sebuah ekstremum lokal
pada himpunan aras $\{g = 0\}$, dengan $\nabla g(a) \neq 0$. Buktikan bahwa
$\nabla f(a) = \lambda\nabla g(a)$ untuk suatu $\lambda$.
\emph{(Parameterkan himpunan arasnya di dekat $a$ lewat teorema fungsi
implisit lalu turunkan $t \mapsto f(\gamma(t))$.)} Penerapannya:
ekstremum $f(x,y) = xy$ pada lingkaran $x^2 + y^2 = 1$.
\end{exercise}

\begin{exercise}[$\star\star$]\label{exo:b2:diffcalc:7}
Tentukan ekstremum $f(x, y) = x^2 + y^2 - xy + x - y$ pada
$\R^2$, lalu maksimum dan minimumnya pada segitiga tertutup bertitik
sudut $(0,0)$, $(1,0)$, $(0,1)$ \emph{(yakni titik kritis dalamnya,
lalu ketiga rusuknya, lalu titik sudutnya)}.
\end{exercise}

\begin{exercise}[$\star\star\star$]\label{exo:b2:diffcalc:8}
Misalkan $f \colon \R^2 \to \R^2$ dengan $f(x, y) = (x + y^2,\; y + x^2)$.
Tunjukkan bahwa $f$ merupakan difeomorfisma lokal di dekat $0$, hitunglah
$\dd(f^{-1})_{(0,0)}$, lalu carilah $r$ terbesar yang membuat $\dd f$
terbalikkan pada bola $\norm{(x,y)}_2 < r$ \emph{(hitunglah
$\det J_f$)}.
\end{exercise}

\begin{exercise}[$\star\star\star$]\label{exo:b2:diffcalc:9}
(Rolle gagal, nilai rerata bertahan) Berikanlah $f \colon \R \to \R^2$
yang $C^1$, dengan $f(0) = f(2\pi)$ tetapi $f'(t) \neq 0$ untuk setiap $t$
(jadi tak ada Rolle yang bernilai vektor). Lalu periksalah pada contohmu
\emph{ketaksamaan} nilai rerata pada \cref{thm:b2:diffcalc:mvi}.
\end{exercise}

\begin{exercise}[$\star$]\label{exo:b2:diffcalc:10}
Hitunglah \omterm{def:b2:diffcalc:differential}{diferensial} dan gradien $f(x) = \norm
x_2^2$ dan $g(x) = \langle Ax, x\rangle$ pada $\R^n$ (dengan $A$ matriks
bujur sangkar yang tak diandaikan \omterm{def:b2:quadratic:adjoint}{simetrik}), beserta Hessian masing-masingnya.
Untuk $A$ yang mana $g$ bersifat cembung?
\end{exercise}

\begin{exercise}[$\star\star$]\label{exo:b2:diffcalc:11}
Misalkan $f \colon \R^n \to \R$ bersifat $C^2$. Buktikan bahwa $f$ cembung bila
dan hanya bila Hessiannya $H_x$ semidefinit positif di setiap
$x$ \emph{(reduksikan ke satu peubah: $t \mapsto f(a + t(b-a))$;
lalu pakailah Taylor--Lagrange pada satu arah, dan untuk konversnya
nilailah $\varphi''$)}.
\end{exercise}

\begin{exercise}[$\star\star\star$]\label{exo:b2:diffcalc:12}
(Sebuah teorema balikan global) Misalkan $g \colon \R^n \to \R^n$ bersifat $C^1$
dengan $\vertiii{\dd g_x} \leq k < 1$ untuk setiap $x$, dan $f =
\mathrm{id} + g$.
\begin{enumerate}
  \item Tunjukkan $\norm{f(x) - f(y)} \geq (1 - k)\norm{x - y}$: jadi $f$
        bersifat injektif, dengan balikan yang \omterm{def:b2:metric:continuity}{kontinu} pada petanya.
  \item Tunjukkan bahwa untuk tiap $y \in \R^n$ pemetaan $x \mapsto y -
        g(x)$ merupakan kontraksi pada ruang \omterm{def:b2:metric:complete}{lengkap} $\R^n$,
        lalu simpulkan lewat teorema titik tetap Banach
        (\cref{thm:b2:metric:banach}) bahwa $f$ bersifat surjektif.
  \item Simpulkan bahwa $f$ merupakan bijeksi $\R^n$ dengan
        balikan \omterm{def:b2:metric:continuity}{Lipschitz} berkonstanta $(1-k)^{-1}$ --- yakni padanan
        \emph{global} bagi \cref{thm:b2:diffcalc:inverse} (yang,
        sebaliknya, bersifat lokal belaka).
\end{enumerate}
\end{exercise}

\section{Soal: Kalkulus matriks --- Jacobi,
eksponensial, dan grup ortogonal}

\begin{problem}\label{pb:b2:diffcalc:1}
Gelanggang terbersih bagi kalkulus \omterm{def:b2:diffcalc:differential}{diferensial} adalah ruang
$\mathcal M_n(\R) \simeq \R^{n^2}$ itu sendiri: sebab pemetaannya yang paling alami
--- hasil kali, balikan, \omterm{def:b2:linalg:det}{determinan}, eksponensial --- punya
\omterm{def:b2:diffcalc:differential}{diferensial} yang keanggunannya memukau. Soal ini menghitung
semuanya: \emph{deret Neumann}, \omterm{def:b2:diffcalc:differential}{diferensial} balikannya,
\emph{rumus Jacobi}\index{rumus Jacobi} bagi
\omterm{def:b2:linalg:det}{determinannya} dengan \emph{rumus Liouville} sebagai panennya,
\emph{eksponensial matriks}\index{eksponensial matriks} dengan
$\det\eu^A = \eu^{\operatorname{tr}A}$, dan akhirnya
grup ortogonal $O_n$ sebagai himpunan aras yang mulus dengan
matriks antisimetrik sebagai ruang singgungnya --- yakni geometri
\omterm{def:b2:diffcalc:differential}{diferensial} dalam embrio. Di sepanjang soal ini, $\vertiii\cdot$ adalah \omterm{thm:b2:nvs:continuouslinear}{norma
operator} yang bawahan terhadap $\norm\cdot_2$, dan $\langle X, Y\rangle =
\operatorname{tr}(X^{\mathsf T}Y)$ hasil kali dalam Frobeniusnya.

\medskip
\noindent\textbf{Bagian I --- Deret Neumann.}
\begin{enumerate}
  \item Buktikan sifat submultiplikatifnya, $\vertiii{AB} \leq
        \vertiii A\,\vertiii B$, lalu turunkan bahwa pemetaan
        polinomial atas $A$ (yakni hasil kali matriks, \omterm{def:b2:linalg:det}{determinan}, trace)
        bersifat \omterm{def:b2:metric:continuity}{kontinu} pada $\mathcal M_n(\R)$.
  \item Untuk $\vertiii X < 1$, tunjukkan bahwa $\sum_{k\geq0}X^k$
        konvergen \omterm{def:b2:series:def}{mutlak} di $\mathcal M_n(\R)$
        (\cref{thm:b2:nvs:absoluteconvergence}), bahwa jumlahnya
        adalah $(I - X)^{-1}$, dan bahwa
        \[
        \vertiii{(I - X)^{-1}} \leq \frac{1}{1 - \vertiii X},
        \qquad
        (I - X)^{-1} = I + X + O\bigl(\vertiii X^2\bigr) .
        \]
  \item Turunkan bahwa $GL_n(\R)$ bersifat \emph{\omterm{def:b2:metric:topology}{terbuka}}: sebab bila $A$
        terbalikkan dan $\vertiii H <
        \frac{1}{\vertiii{A^{-1}}}$, maka $A + H$ bersifat
        terbalikkan. Turunkan pula bahwa $GL_n(\R)$ bersifat \emph{padat}
        di $\mathcal M_n(\R)$ \emph{(usiklah $A$ dengan
        $\varepsilon I$: sebab $\det(A + \varepsilon I)$ merupakan polinomial
        tak nol dalam $\varepsilon$)}.
  \item Tunjukkan bahwa pemetaan pembalikan $\Phi(A) = A^{-1}$ bersifat
        \omterm{def:b2:metric:continuity}{kontinu} pada $GL_n(\R)$.
\end{enumerate}

\medskip
\noindent\textbf{Bagian II --- \omterm{def:b2:diffcalc:differential}{Diferensial} pertama.}
\begin{enumerate}[resume]
  \item Tunjukkan bahwa pemetaan pengkuadratan $A \mapsto A^2$ bersifat
        \omterm{def:b2:diffcalc:differential}{terdiferensialkan} dengan \omterm{def:b2:diffcalc:differential}{diferensial} $H \mapsto AH + HA$,
        dan lebih umum bahwa $A \mapsto A^k$ punya
        \omterm{def:b2:diffcalc:differential}{diferensial} $H \mapsto \sum_{i=0}^{k-1}
        A^iHA^{k-1-i}$. Mengapa kita \emph{tak} boleh menulis $kA^{k-1}H$
        secara umum?
  \item Buktikan bahwa $\Phi(A) = A^{-1}$ \omterm{def:b2:diffcalc:differential}{terdiferensialkan} pada
        $GL_n(\R)$ dengan
        \[
        \dd\Phi_A(H) = -A^{-1}HA^{-1}
        \]
        \emph{(tulislah $(A + H)^{-1} = (I + A^{-1}H)^{-1}A^{-1}$
        lalu uraikan lewat pertanyaan 2)}. Periksalah rumusnya terhadap
        kasus skalar $n = 1$.
  \item Untuk kurva $C^1$ $t \mapsto A(t) \in GL_n(\R)$, turunkan
        $\bigl(A(t)^{-1}\bigr)' = -A^{-1}A'A^{-1}$, lalu uraikan
        $t \mapsto (I + tB)^{-1}$ sampai orde pertama di $t = 0$.
  \item Hitunglah \omterm{def:b2:diffcalc:differential}{diferensial} $f(A) =
        \operatorname{tr}(A^k)$ lalu kenali gradiennya bagi
        hasil kali dalam Frobeniusnya:
        \[
        \dd f_A(H) = k\operatorname{tr}\bigl(A^{k-1}H\bigr),
        \qquad
        \nabla f(A) = k\,\bigl(A^{k-1}\bigr)^{\mathsf T} .
        \]
  \item Pertanyaan yang sama untuk $f(A) = \operatorname{tr}
        (A^{\mathsf T}A) = \norm A_F^2$: yakni \omterm{def:b2:diffcalc:differential}{diferensialnya},
        gradiennya, dan Hessiannya (yang tetap); lalu simpulkan bahwa
        $\norm\cdot_F^2$ bersifat cembung tegas.
\end{enumerate}

\medskip
\noindent\textbf{Bagian III --- Rumus Jacobi.}
\begin{enumerate}[resume]
  \item Buktikan
        \[
        \det(I + H) = 1 + \operatorname{tr}H +
        O\bigl(\vertiii H^2\bigr)
        \]
        \emph{(uraikan $\det(e_1 + h_1, \dots, e_n + h_n)$ lewat
        kemultilinearan pada kolomnya: sebab suku dengan sekurangnya dua
        kolom-$h$ bernilai $O(\vertiii H^2)$)}: jadi $\dd(\det)_I =
        \operatorname{tr}$.
  \item Untuk $A$ yang terbalikkan, turunkan
        \[
        \dd(\det)_A(H) = \det(A)\,
        \operatorname{tr}\bigl(A^{-1}H\bigr) .
        \]
  \item Tunjukkan bahwa untuk \emph{setiap} $A$ (terbalikkan atau tidak),
        berlaku $\frac{\partial\det}{\partial a_{ij}}(A) = C_{ij}$, yakni
        kofaktor $(i,j)$-nya \emph{(lewat penguraian Laplace sepanjang baris
        $i$)}, sehingga dengan adjugatnya $\operatorname{adj}A =
        \operatorname{com}(A)^{\mathsf T}$:
        \[
        \dd(\det)_A(H) =
        \operatorname{tr}\bigl(\operatorname{adj}(A)\,H\bigr),
        \qquad
        \nabla(\det)(A) = \operatorname{com}(A) ,
        \]
        yang memulihkan pertanyaan 11 ketika $A$ terbalikkan
        (sebab $\operatorname{adj}A = \det(A)A^{-1}$). Inilah
        \emph{rumus Jacobi}: $\bigl(\det
        A(t)\bigr)' = \operatorname{tr}\bigl(
        \operatorname{adj}(A(t))\,A'(t)\bigr)$.
  \item (Rumus Liouville) Misalkan $A(t)$ kurva matriks $C^1$
        yang memenuhi persamaan \omterm{def:b2:diffcalc:differential}{diferensial} linear
        $A'(t) = M(t)A(t)$. Buktikan
        \[
        \bigl(\det A(t)\bigr)' =
        \operatorname{tr}\bigl(M(t)\bigr)\,\det A(t),
        \qquad\text{sehingga}\qquad
        \det A(t) = \det A(0)\,
        \exp\Bigl(\int_0^t\operatorname{tr}M\Bigr)
        \]
        \emph{(pakailah $\operatorname{adj}(A)\,A = \det(A)I$ dan
        keinvarianan \omterm{def:b2:structures:generated}{siklik} tracenya)} --- yakni kesamaan Wronskian
        yang kelak dipakai bab persamaan \omterm{def:b2:diffcalc:differential}{diferensial} secara terus-menerus.
  \item Tunjukkan bahwa $SL_n(\R) = \{\det = 1\}$ merupakan himpunan aras
        yang mulus: sebab di setiap $A \in SL_n(\R)$ \omterm{def:b2:diffcalc:differential}{diferensial}
        $\dd(\det)_A$ merupakan pemetaan linear yang \emph{surjektif} ke
        $\R$ \emph{(nilailah ia di $H = \frac1nA$)}.
\end{enumerate}

\medskip
\noindent\textbf{Bagian IV --- \omterm{ex:b2:nvs:matrixexp}{Eksponensial matriks}.}
\begin{enumerate}[resume]
  \item Tunjukkan bahwa $\eu^A = \sum_{k\geq0}\frac{A^k}{k!}$
        konvergen \omterm{def:b2:series:def}{mutlak} untuk setiap $A$, dan normal pada setiap
        bola, dengan $\vertiii{\eu^A} \leq
        \eu^{\vertiii A}$; dan bahwa $\eu^A$ bergantung
        secara \omterm{def:b2:metric:continuity}{kontinu} pada $A$.
  \item Buktikan bahwa $AB = BA$ mengakibatkan $\eu^{A+B} =
        \eu^A\eu^B$ \emph{(lewat hasil kali Cauchy, yang sah berkat
        kekonvergenan \omterm{def:b2:series:def}{mutlaknya})}; lalu turunkan bahwa $\eu^A$ selalu
        terbalikkan, dengan balikan $\eu^{-A}$: jadi $\exp$ memetakan
        $\mathcal M_n(\R)$ ke dalam $GL_n(\R)$.
  \item Tunjukkan bahwa $t \mapsto \eu^{tA}$ bersifat $C^1$ (bahkan
        $C^\infty$) dengan
        \[
        \frac{\dd}{\dd t}\,\eu^{tA} = A\,\eu^{tA} =
        \eu^{tA}A
        \]
        \emph{(turunkan deretnya suku demi suku pada ruas:
        sebab deret turunannya konvergen normal)}.
  \item Buktikan kesamaan
        \[
        \det\bigl(\eu^{A}\bigr) = \eu^{\operatorname{tr}A}
        \]
        \emph{(terapkan rumus Liouville pada pertanyaan 13 untuk
        $A(t) = \eu^{tA}$)}. Pemeriksaan kewarasannya: kasus $n = 1$; lalu
        $A$ yang nilpoten; dan matriks bertrace nol yang mendarat di
        $SL_n(\R)$.
  \item Tunjukkan $\eu^H = I + H + O(\vertiii H^2)$, jadi $\exp$ bersifat
        \omterm{def:b2:diffcalc:differential}{terdiferensialkan} di $0$ dengan $\dd(\exp)_0 =
        \mathrm{id}$; lalu simpulkan lewat teorema fungsi
        balikan (\cref{thm:b2:diffcalc:inverse}) bahwa $\exp$
        merupakan difeomorfisma $C^1$ dari sebuah persekitaran $0$
        ke sebuah persekitaran $I$: jadi setiap matriks yang dekat ke
        identitasnya punya logaritma.
  \item Tunjukkan bahwa $\exp$ memetakan matriks \omterm{def:b2:quadratic:adjoint}{simetrik} ke matriks
        \omterm{def:b2:quadratic:adjoint}{simetrik} yang \emph{definit positif}, secara bijektif
        \emph{(diagonalkan; sebab balikannya adalah logaritma
        spektralnya)}.
\end{enumerate}

\medskip
\noindent\textbf{Bagian V --- Grup ortogonal sebagai sebuah himpunan
aras.}
\begin{enumerate}[resume]
  \item Misalkan $F(A) = A^{\mathsf T}A$, dari $\mathcal M_n(\R)$ ke
        matriks \omterm{def:b2:quadratic:adjoint}{simetrik} $S_n$. Hitunglah $\dd F_A(H) =
        A^{\mathsf T}H + H^{\mathsf T}A$ lalu tunjukkan bahwa di setiap
        $A \in O_n = F^{-1}(I)$ \omterm{def:b2:diffcalc:differential}{diferensial} ini bersifat
        \emph{surjektif} ke $S_n$ \emph{(diberikan $S \in S_n$,
        cobalah $H = \frac12 AS$)}: jadi $O_n$ merupakan himpunan aras yang mulus,
        berdimensi $n^2 - \frac{n(n+1)}2 = \frac{n(n-1)}2$.
  \item Tunjukkan bahwa setiap kurva $C^1$ $A(t) \in O_n$ dengan $A(0) =
        I$ punya kecepatan $A'(0)$ yang \emph{antisimetrik}, dan
        sebaliknya bahwa untuk $K$ yang antisimetrik kurva
        $\eu^{tK}$ tinggal di $O_n$: jadi ruang singgung $O_n$
        di $I$ persis merupakan matriks antisimetriknya.
  \item Tunjukkan $\det\eu^{K} = 1$ untuk $K$ yang antisimetrik (pertanyaan
        18): jadi kurva eksponensialnya hidup di grup perputaran
        $SO_n$. Hitunglah ia selengkapnya untuk $n = 2$: dengan $J =
        \begin{pmatrix}0 & -1\\ 1 & 0\end{pmatrix}$, buktikanlah
        \[
        \eu^{\theta J} = \begin{pmatrix} \cos\theta &
        -\sin\theta\\ \sin\theta & \cos\theta\end{pmatrix} :
        \]
        jadi eksponensial matriksnya \emph{adalah} perputaran sebesar
        $\theta$, dan definisi deret bagi kosinus dan sinus
        muncul kembali di dalam sebuah matriks.
  \item (Siasat kepadatan) Dengan memakai kepadatan $GL_n(\R)$
        (pertanyaan 3) dan kekontinuannya, perluaslah dari yang terbalikkan ke
        semua matriks kesamaan
        \[
        \operatorname{adj}(AB) =
        \operatorname{adj}(B)\operatorname{adj}(A)
        \]
        \emph{(sebab untuk $A, B$ yang terbalikkan kedua ruasnya sama dengan
        $\det(AB)(AB)^{-1}$; dan kedua ruasnya polinomial dalam
        entrinya)}.
  \item Rangkuman. Satu kalimat untuk masing-masing: (i) bab
        terdahulu mana yang menyediakan mesin tiap Bagiannya
        (yakni kelengkapan dan aljabar bernorma; teorema
        spektral; teorema fungsi balikan); (ii) rumus mana
        pada soal ini yang kelak disandari bab persamaan
        \omterm{def:b2:diffcalc:differential}{diferensial}, dan di mana; (iii) apa yang dikatakan
        $\dd(\det)_I = \operatorname{tr}$ dan $\det\eu^A =
        \eu^{\operatorname{tr}A}$ tentang trace dan
        \omterm{def:b2:linalg:det}{determinan} sebagai ``volume infinitesimal dan global'';
        (iv) apa yang dijadikan jilid Tahun ke-3 dari pertanyaan
        21--23 (yakni grup Lie beserta aljabar Lie-nya).

\end{enumerate}
\end{problem}
